%%%%%%%%%%%%%%%%%%%%%%%%%%%%%%%%%%%%%%%%%%%%%%%%%%%%%%%%%%%%%%%%%%
\section{Discussion}\label{sec:discussion}
%%%%%%%%%%%%%%%%%%%%%%%%%%%%%%%%%%%%%%%%%%%%%%%%%%%%%%%%%%%%%%%%%%

\todo{Not drafted in this pass. Intended themes, to be written once the essential evidence in paper-data/manuscript\_todos.md is available:
\begin{itemize}
	\item what an FSI verification claim should contain (reference category, levels, QoIs, iterative-error evidence), and how this differs from current practice in the literature;
	\item the value of exact references for exposing and tracing sub-nominal orders (womersleyTube: an $O(\Delta t)$ inconsistency in the truncated-domain boundary treatment, not the coupling) and for separating coupling, solid and temporal errors (ringAddedMass);
	\item limits of the present suite: no compressible fluids, no contact, no free surfaces, limited non-Newtonian/hyperelastic coverage, a single code family (independent references only for collapsibleChannel, FSI1, and pending COMSOL cases);
	\item model-form differences in code-to-code comparisons (beam versus continuum in collapsibleChannel; Hookean versus St Venant--Kirchhoff in 3dTube);
	\item verification versus validation: where experimental data (Hessenthaler) fit, and why they are kept separate;
	\item recommendations to benchmark authors (tabulate values, state mesh/time studies, provide machine-readable references and source).
\end{itemize}}
